\begin{center}\begin{tikzpicture}[font=\small,>=stealth,line join=round]
\draw[gray!40,rounded corners=3pt] (0,0) rectangle (14.2,8.3);\node at (7.1,7.700000000000001) {同样得到偶函数，外层的性质可以完全不同};
\begin{scope}[shift={(4.1,2.068181818181818)},x=2.0cm,y=0.41818181818181815cm]\draw[->,gray] (-1.5,0)--(1.5,0) node[right]{$x$};\draw[->,gray] (0,-1)--(0,10) node[above]{$y$};\draw[blue,thick,solid,domain=-1.2:1.2,samples=180] plot(\x,{(\x)^2+1});\draw[violet,thick,dashed,domain=-1.2:1.2,samples=180] plot(\x,{((\x)^2+1)^2+(\x)^2+1});\end{scope}\node[anchor=west,align=left,text width=5.0cm] at (8.4,4.0) {内层都为 $g=x^2+1$\\[10pt]蓝：外层 $f(u)=u$（奇）\\[10pt]紫：$f(u)=u^2+u$（非奇非偶）\\[10pt]两个整体都是偶函数};
\node[text width=12.6cm,align=center] at (7.1,.65) {两条实际曲线说明：仅从整体为偶，无法反推外层的奇偶性。};
\end{tikzpicture}\end{center}
